% Working proof fragment, not integrated into a release.
% Depends on the separately proved Lemma lem:gn-norm-transfer.
% No Frei--Tinkova theorem is assumed in this fragment.
% Their ideal-norm Type I/II strategy is retained as methodological attribution.
\subsection{Gaussian ideal Vaughan decomposition and Type I/II estimates}
\label{gi-section}

We give the ideal decomposition and the subsequent integer-parameter
recombination in full.  The architecture is that of
\cite[Section~3]{GTMobius2012}, with the Gaussian ideal-norm adaptation
attributed to Frei--Tinkova.  The analytic inputs used here are the
separately proved Lemma~\ref{lem:gn-norm-transfer} and the published
one-parameter results \cite[Theorem~2.9 and Lemma~4.5]{GTOrbits2012}.
In particular, no theorem of Frei--Tinkova is an assumed input.

Fix the Gaussian ray modulus $H_0$.  All ideals in this subsection are
nonzero integral ideals of $\mathbb Z[i]$.  Write $\mathrm N\mathfrak a$
for the absolute ideal norm, $\mu_E$ for ideal Mobius, and
$\Lambda_E(\mathfrak p^j)=\log\mathrm N\mathfrak p$ for $j\geq1$, with
$\Lambda_E$ zero on other ideals.  A ray character $\chi$ is extended by
zero on ideals not coprime to $H_0$.  It is completely multiplicative
and satisfies $|\chi|\leq1$.  We use a Lipschitz norm that includes the
supremum norm.  All nilmanifolds have connected, simply connected
nilpotent covering group, a cocompact lattice, and an adapted rational
Mal'cev basis, with the filtration and polynomial-sequence conventions
of \cite{GTOrbits2012}.

\begin{lemma}[Ideal multiplicities and moments]\label{gi-moments}
Let $a_E(n)=\#\{\mathfrak a:\mathrm N\mathfrak a=n\}$.  For $X\geq1$,
\begin{align}
 a_E(n)&\leq\tau(n),\label{gi-multiplicity}\\
 \sum_{n\leq X}a_E(n)^2&\leq X(1+\log X)^3,&
 \sum_{n\leq X}\frac{a_E(n)^2}{n}&\leq(1+\log X)^4.\label{gi-two-moments}
\end{align}
Moreover
\begin{equation}\label{gi-ideal-count}
 A_E(X):=\#\{\mathfrak a:\mathrm N\mathfrak a\leq X\}
       =\frac{\pi}{4}X+O(\sqrt X+1).
\end{equation}
In particular $A_E(X)\ll X$, and the number of ideals of norm in
$(B/2,B]$ is comparable to $B$ for all sufficiently large $B$.
\end{lemma}
\begin{proof}
At a split rational prime $p$ one has $a_E(p^k)=k+1$; at an inert prime
it is one for even $k$ and zero for odd $k$; and $a_E(2^k)=1$.
Multiplicativity proves \eqref{gi-multiplicity}.  Prime by prime,
$(k+1)^2\leq\binom{k+3}{3}$, so $\tau(n)^2\leq d_4(n)$.
Summing over ordered quadruples with product at most $X$, and summing
the last coordinate first, bounds $\sum_{n\leq X}d_4(n)$ by
$X(1+\log X)^3$.  Bounding each of the four harmonic sums separately
gives $\sum_{n\leq X}d_4(n)/n\leq(1+\log X)^4$.
Every ideal has exactly four Gaussian generators.  Counting lattice
points in the disk of radius $\sqrt X$, and dividing by four, proves
\eqref{gi-ideal-count}: the unit squares meeting its boundary lie
between circles of radii $\sqrt X-\sqrt2$ and $\sqrt X+\sqrt2$ and
number $O(\sqrt X+1)$.
\end{proof}

\begin{lemma}[Exact ideal Vaughan identity]\label{gi-vaughan}
Let $N\geq8$, $U=V=N^{1/3}$, $L=\log(2N)$, and let
$f:\mathbb N\longrightarrow\mathbb C$ satisfy $|f|\leq1$ and vanish
outside $(N/2,N]$.  Then
\begin{equation}\label{gi-vaughan-sum}
 \sum_{\mathfrak a}\Lambda_E(\mathfrak a)\chi(\mathfrak a)
                       f(\mathrm N\mathfrak a)=T_0-T_1+T_2,
\end{equation}
where
\begin{align*}
 T_0&=\sum_{\mathrm N\mathfrak d\leq U}\mu_E(\mathfrak d)\chi(\mathfrak d)
       \sum_{\mathfrak b}\log(\mathrm N\mathfrak b)\chi(\mathfrak b)
                  f(\mathrm N\mathfrak d\,\mathrm N\mathfrak b),\\
 T_1&=\sum_{\mathrm N\mathfrak d\leq UV}\alpha_{\mathfrak d}\chi(\mathfrak d)
       \sum_{\mathfrak b}\chi(\mathfrak b)
                  f(\mathrm N\mathfrak d\,\mathrm N\mathfrak b),\\
 T_2&=\sum_{V<\mathrm N\mathfrak d\leq N/U}\beta_{\mathfrak d}\chi(\mathfrak d)
       \sum_{\mathrm N\mathfrak b>U}\mu_E(\mathfrak b)\chi(\mathfrak b)
                  f(\mathrm N\mathfrak d\,\mathrm N\mathfrak b),
\end{align*}
with
\begin{align*}
 \alpha_{\mathfrak d}
   &=\sum_{\substack{\mathfrak b\mathfrak c=\mathfrak d\\
              \mathrm N\mathfrak b\leq U,\ \mathrm N\mathfrak c\leq V}}
                        \mu_E(\mathfrak b)\Lambda_E(\mathfrak c),&
 \beta_{\mathfrak d}
   &=\sum_{\substack{\mathfrak c\mid\mathfrak d\\\mathrm N\mathfrak c>V}}
                        \Lambda_E(\mathfrak c).
\end{align*}
Both coefficients have absolute value at most $L$.
\end{lemma}
\begin{proof}
Convolution on ideals satisfies
$\mu_E*1=\mathbf1_{\mathfrak a=(1)}$, $1*\Lambda_E=\log\mathrm N$,
and $\mu_E*\log\mathrm N=\Lambda_E$, as one checks on prime powers.
Splitting $\mu_E$ at $U$ and $\Lambda_E$ at $V$ gives the exact identity
\begin{equation}\label{gi-convolution-identity}
 \Lambda_E=\mu_{\leq U}*\log\mathrm N+\Lambda_{\leq V}
       -\mu_{\leq U}*1*\Lambda_{\leq V}
       +\mu_{>U}*1*\Lambda_{>V}.
\end{equation}
The second term vanishes on the support of $f$.  Complete multiplicativity
of $\chi$, including its zero extension, permits separation of each
convolution factor and yields \eqref{gi-vaughan-sum} with exactly the
stated cutoffs.  Finally
$\sum_{\mathfrak c\mid\mathfrak d}\Lambda_E(\mathfrak c)
=\log\mathrm N\mathfrak d$ and every summand is nonnegative.  This
proves the coefficient bounds, with no discarded endpoint or error term.
\end{proof}

\begin{proposition}[Quantitative extraction]\label{gi-extraction}
There are absolute constants $c_0>0$ and $N_0$ such that the following
holds.  Under Lemma~\ref{gi-vaughan}, suppose $N\geq N_0$,
$0<\epsilon\leq1$, and
\begin{equation}\label{gi-large-correlation}
 \left|\sum_{\mathfrak a}\Lambda_E(\mathfrak a)\chi(\mathfrak a)
                         f(\mathrm N\mathfrak a)\right|\geq\epsilon NL.
\end{equation}
Put $\theta=c_0\epsilon^4/L^{24}$.  At least one of the following holds.
\begin{enumerate}
\item Some dyadic $R\leq2N^{1/3}$ has at least $\theta R$ integers
$r\in(R/2,R]$ satisfying
\[
 \left|\sum_{\mathfrak b}\log(\mathrm N\mathfrak b)\chi(\mathfrak b)
                   f(r\mathrm N\mathfrak b)\right|\geq\theta N/r.
\]
\item Some dyadic $R\leq2N^{2/3}$ has at least $\theta R$ integers
$r\in(R/2,R]$ satisfying the same inequality without the logarithm.
\item There are dyadic $R,B$ with
\begin{equation}\label{gi-dyadic-ranges}
 N^{1/3}<R\leq2N^{2/3},\qquad N/2<RB<4N,
\end{equation}
and at least $\theta R^2$ pairs $r,r'\in(R/2,R]$ satisfying
\begin{equation}\label{gi-type-two-correlation}
 \left|\sum_{B/2<\mathrm N\mathfrak b\leq B}
       f(r\mathrm N\mathfrak b)\overline{f(r'\mathrm N\mathfrak b)}
       \right|\geq\theta B.
\end{equation}
\end{enumerate}
\end{proposition}
\begin{proof}
One of $|T_i|$ is at least $\epsilon NL/3$.  For $T_0$ put
\[
 G(r)=\frac{r}{NL}\left|\sum_{\mathfrak b}\log(\mathrm N\mathfrak b)
                   \chi(\mathfrak b)f(r\mathrm N\mathfrak b)\right|;
\]
for $T_1$ use $r/N$ times the absolute unweighted inner sum.  In either
case Lemma~\ref{gi-moments} gives $G\ll1$ and
\[
 \epsilon\ll\sum_{r\leq N^e}\frac{a_E(r)G(r)}r,
 \qquad e=1/3\text{ or }2/3.
\]
Cauchy--Schwarz and \eqref{gi-two-moments} imply
$\sum_{r\leq N^e}G(r)^2/r\gg\epsilon^2/L^4$.
There are $O(L)$ dyadic cells, so one satisfies
\[
 \sum_{R/2<r\leq R}G(r)^2\gg R\epsilon^2/L^5.
\]
Since $G\ll1$, at least $cR\epsilon^2/L^5$ terms have
$G(r)\geq c\epsilon/L^{5/2}$.  Both the count and the claimed amplitudes
follow after decreasing $c_0$.

For $T_2$, split both ideal norms into dyadic cells.  There are $O(L^2)$
rectangles, and a nonzero cell satisfies \eqref{gi-dyadic-ranges} and
$B\geq N^{1/3}/4$.  A rectangle contributes at least $c\epsilon N/L$;
we use the weaker bound $c\epsilon N/L^2$.  Grouping $\mathfrak d$ by
$r=\mathrm N\mathfrak d$ writes this contribution as
\[
 T=\sum_{R/2<r\leq R}A_r\sum_{\mathfrak b\in\mathcal B}c_{\mathfrak b}
                   f(r\mathrm N\mathfrak b),\qquad
 \mathcal B=\{\mathfrak b:B/2<\mathrm N\mathfrak b\leq B\}.
\]
Here $|A_r|\leq L a_E(r)$ and $|c_{\mathfrak b}|\leq1$; the restriction
$\mathrm N\mathfrak b>U$ belongs to $c_{\mathfrak b}$.  All product
cutoffs remain in $f$, so the domain is rectangular.  By
\eqref{gi-two-moments}, $\sum_r|A_r|^2\ll RL^5$.  Thus, with
$F_{r,\mathfrak b}=f(r\mathrm N\mathfrak b)$,
\begin{equation}\label{gi-first-cauchy}
 T_*:=\sum_r\left|\sum_{\mathfrak b\in\mathcal B}
                          c_{\mathfrak b}F_{r,\mathfrak b}\right|^2
          \gg \epsilon^2 RB^2/L^9.
\end{equation}
Put $M=|\mathcal B|\asymp B$.  Expansion and Cauchy--Schwarz in the
ordered pair $(\mathfrak b,\mathfrak b')$ give
\begin{align}
 T_*^2
 &=\left|\sum_{\mathfrak b,\mathfrak b'}
 c_{\mathfrak b}\overline{c_{\mathfrak b'}}
 \sum_r F_{r,\mathfrak b}\overline{F_{r,\mathfrak b'}}\right|^2\notag\\
 &\leq M^2\sum_{\mathfrak b,\mathfrak b'}
       \left|\sum_rF_{r,\mathfrak b}\overline{F_{r,\mathfrak b'}}\right|^2
 =M^2\sum_{r,r'}\left|\sum_{\mathfrak b}
                  F_{r,\mathfrak b}\overline{F_{r',\mathfrak b}}\right|^2.
 \label{gi-second-cauchy}
\end{align}
Consequently
\[
 \mathbb E_{r,r'}\left|\mathbb E_{\mathfrak b\in\mathcal B}
              F_{r,\mathfrak b}\overline{F_{r',\mathfrak b}}\right|^2
       \gg\epsilon^4/L^{18}.
\]
Every summand is at most one.  Thresholding gives
\eqref{gi-type-two-correlation} and its stated count, with the weaker
common parameter $\theta=c_0\epsilon^4/L^{24}$.  The character and Mobius
coefficients have disappeared.  All conjugates in
\eqref{gi-second-cauchy} are necessary.
\end{proof}

\paragraph{Progressions and interval restrictions.}\label{gi-progressions}
In the application,
\[
 f(n)=1_{N/2<n\leq N}\,1_{u<n\leq v}\,1_{n\equiv b\pmod q}
                              F(g(n)\Gamma).
\]
For Type I the congruence $rn\equiv b\pmod q$ is either insoluble or
one class modulo $q/\gcd(r,q)$.  In Type II its intersection with
$r'n\equiv b\pmod q$ is either empty or one class modulo
\[
 \operatorname{lcm}\bigl(q/\gcd(r,q),q/\gcd(r',q)\bigr),
\]
which divides $q$.  The remaining restrictions are interval intersections.
Lemma~\ref{lem:gn-norm-transfer} covers them without a lower bound on
their lengths.  On the product nilmanifold the function is
$F(x)\overline{F(y)}$, whose Haar integral is
$|\int F|^2=0$.  After normalizing $\|F\|_{\mathrm{Lip}}=1$, its
Lipschitz norm for the sum metric is at most two.  Thus all mean-zero,
modulus and interval hypotheses for product norm transfer are satisfied.

\begin{lemma}[Integer-parameter recombination]\label{gi-recombination}
For each $m,d$ there are $C_{\rm rec},K_{\rm rec}\geq2$ with the following
property.  Let $0<\eta<\eta_0$, $N\geq\eta^{-K_{\rm rec}}$, and let $g$
have degree at most $d$ on an $m$-dimensional nilmanifold with an
$\eta^{-1}$-rational adapted basis.  Each of the following implies that
$(g(n)\Gamma)_{n\leq N}$ is not totally $\eta^{C_{\rm rec}}$-equidistributed.
\begin{enumerate}
\item For at least $\eta R$ values $r\in(R/2,R]$, where
$R\leq2N^{2/3}$, the orbit $(g(rn)\Gamma)_{n\leq\lceil N/r\rceil}$
is not totally $\eta$-equidistributed.
\item For at least $\eta R^2$ pairs $r,r'\in(R/2,R]$, where
$N^{1/3}/4\leq R\leq2N^{2/3}$ and $N/2<RB<4N$, the product orbit
$(g(rn),g(r'n))_{n\leq B}$ is not totally $\eta$-equidistributed.
\end{enumerate}
\end{lemma}
\begin{proof}
All constants in this proof depend only on $m,d$.  Failure of total
equidistribution gives a progression $a+q[\ell]\subseteq[M]$ with
$\ell\geq\eta M$, $q\leq2/\eta$, and $|a|\leq2M$ on which ordinary
$\eta$-equidistribution fails.  Apply
\cite[Theorem~2.9]{GTOrbits2012}, in dimension $m$ or $2m$.  A nonzero
horizontal character has magnitude and restricted smoothness norm at
most $\eta^{-C_1}$, after shrinking $\eta_0$ to absorb fixed factors.

Here is the required removal of a possibly large shift.  Write the scalar
phase $p(a+qt)$ in monomials.  Multiplication by $d!$ clears conversion
from its binomial coefficients.  Inverting $t=(n-a)/q$ and multiplying
by $q^d$ clears every remaining denominator.  The error in the coefficient
of $n^i$ is bounded by
\begin{equation}\label{gi-shift-removal}
 C_d\eta^{-C_1}\sum_{j=i}^d q^{d-j}|a|^{j-i}\ell^{-j}
       \leq\eta^{-C_2}M^{-i},\qquad i\geq1.
\end{equation}
The integer parts disappear modulo one.  Thus a nonzero integer multiple
of the character has magnitude and $C^\infty[M]$ norm at most
$\eta^{-C_2}$.  This argument does not treat $a=O(M)$ as a bounded-height
rational parameter.  There are at most $\eta^{-C_3}$ possible resulting
characters, since their lattice has rank at most $2m$.  Pigeonholing
therefore fixes one character on a set of density at least $\eta^{C_4}$.

In Type I write
$\psi(g(n))=\sum_{j=0}^d\beta_jn^j$ modulo one.  For the selected set
$S\subseteq(R/2,R]$, coefficient conversion gives
\begin{equation}\label{gi-one-coefficients}
 \|d!\beta_jr^j\|_{\mathbb R/\mathbb Z}
       \leq\eta^{-C_5}(R/N)^j,\qquad 1\leq j\leq d.
\end{equation}
If $R\leq\eta^{-C_6}$ for a sufficiently large fixed $C_6$, choose one
$r\in S$ and multiply the character by $d!r^d$.  Multiplication of the
$j$th bound by $r^{d-j}$ gives magnitude and $C^\infty[N]$ norm bounded
by a fixed power of $\eta^{-1}$.  This treats every small dyadic scale,
including $R=1$.

Otherwise apply \cite[Lemma~4.5]{GTOrbits2012} to
$r\mapsto d!\beta_jr^j$ on $[\lceil R\rceil]$.  The selected set has
density at least a fixed power of $\eta$, and its values lie in an arc
of length at most $2\eta^{-C_5}(R/N)^j$.  Taking $K_{\rm rec}$ large,
the bound $R\leq2N^{2/3}$ makes this length at most half the density.
The choice of $C_6$ also places the recurrence length beyond the fixed
inverse density powers.  The published recurrence lemma supplies
nonzero integers $t_j$, $|t_j|\leq\eta^{-C_7}$, such that
\begin{equation}\label{gi-recurrent-coefficients}
 \|t_j(d!)^2\beta_j\|_{\mathbb R/\mathbb Z}
                    \leq\eta^{-C_8}N^{-j}.
\end{equation}
The extra factorial clears the leading-coefficient conversion.  Taking
the product of the $t_j$ and the fixed factorials gives a single nonzero
horizontal character of magnitude and $C^\infty[N]$ norm at most
$\eta^{-C_9}$.

For Type II the pigeonholed character is $\psi_1\oplus\psi_2$, not both
zero.  Write the two scalar phases with coefficients $\beta_j,\beta'_j$.
The information replacing \eqref{gi-one-coefficients} is
\begin{equation}\label{gi-two-coefficients}
 \|d!(\beta_jr^j+\beta'_jr'^j)\|_{\mathbb R/\mathbb Z}
                        \leq\eta^{-C_5}(R/N)^j.
\end{equation}
Interchange the two coordinates if necessary so that $\psi_1\ne0$.
Fix an $r'$ with at least $\eta^{C_4}R$ partners.  Equation
\eqref{gi-two-coefficients} puts $d!\beta_jr^j$ in an arc centered at
$-d!\beta'_jr'^j$.  The same published recurrence lemma permits this
arbitrary center and gives \eqref{gi-recurrent-coefficients} for
$\psi_1$.  Its large-length requirement follows from
$R\geq N^{1/3}/4$ and a sufficiently large $K_{\rm rec}$.  There is no
substitution of $(r-r_0)^j$ for $r^j-r_0^j$.

Finally let $\Psi$ denote the nonzero character just obtained.  Its
binomial expansion shows that for
$1\leq n\leq c\eta^{2C_9+4}N$ the phase $\Psi(g(n))$ stays within
$1/10$ of $\Psi(g(0))$.  Choose a mean-zero Lipschitz circle function
equal to one on that arc and compose it with $\Psi$.  Its mean is zero
because a nonzero horizontal character pushes Haar measure to circle
Haar measure.  Its Lipschitz norm is bounded by a fixed power of
$\eta^{-1}$ by \cite[Appendix~A]{GTOrbits2012}.  Choose $C_{\rm rec}$
larger than the length and Lipschitz exponents, and then enlarge
$K_{\rm rec}$ so the initial interval is nonempty.  Its average one
contradicts total $\eta^{C_{\rm rec}}$-equidistribution.  This proves
both assertions.  All exponent constants are finite combinations of
the published one-parameter constants and degree-$d$ polynomial algebra.
\end{proof}

\begin{proposition}[Power-form Gaussian ideal minor-arc bound]
\label{gi-minor-arc}
For each $m,d$ there is $c_{\rm G}>0$ with the following property.
Fix $H_0$ and a finite-order Gaussian ray character $\chi$ of modulus
$H_0$, extended by zero.  Let $G/\Gamma$ have dimension $m$, filtration
degree $d$, and a $Q$-rational adapted basis, with $Q\geq2$.  Suppose
$(g(n)\Gamma)_{n\leq N}$ is totally $\delta$-equidistributed, where
$0<\delta<1$.  If $F$ is Lipschitz and $\int F=0$, then for every
interval progression $P\subseteq[N]$ of common difference at most $Q$,
\begin{equation}\label{gi-power-bound}
 \left|\sum_{\mathrm N\mathfrak a\leq N}\Lambda_E(\mathfrak a)
       \chi(\mathfrak a)1_P(\mathrm N\mathfrak a)
       F(g(\mathrm N\mathfrak a)\Gamma)\right|
 \ll_{H_0,m,d} QN\log(2N)^7\delta^{c_{\rm G}}\|F\|_{\mathrm{Lip}}.
\end{equation}
No lower bound on the length of $P$ or on $\delta$ is required.
\end{proposition}
\begin{proof}
Normalize $\|F\|_{\mathrm{Lip}}=1$.  Take the smaller positive exponents
$\kappa,\sigma$ in Lemma~\ref{lem:gn-norm-transfer} for dimensions $m$
and $2m$, and decrease $\kappa$ to at most one.  Partial summation gives
the same transfer bound with an extra factor $\log(2M)$ for a
$\log\mathrm N\mathfrak b$ weight: the unweighted lemma holds for
every partial interval, and $\int_1^Mdt/t\leq\log(2M)$.
Let $C_{\rm rec},K_{\rm rec}$ be as in Lemma~\ref{gi-recombination},
and choose
\begin{equation}\label{gi-exponent-choice}
 a_*=\frac1{4C_{\rm rec}},\qquad
 c_* =\frac{a_*\kappa}{32},\qquad
 s_* =\min\left\{\frac1{100},\frac1{4a_*K_{\rm rec}},
                             \frac{\sigma}{12a_*}\right\}.
\end{equation}
First assume $N^{-s_*}\leq\delta\leq\delta_0$, with the fixed cutoff
$\delta_0$ sufficiently small, and restrict the sum to $(N/2,N]$.
We claim that its absolute value is
\begin{equation}\label{gi-shell-bound}
                  \ll QNL^7\delta^{c_*},\qquad L=\log(2N).
\end{equation}
The elementary bound is $O(NL)$.  Thus the claim is immediate when
$QL^6\delta^{c_*}\geq1$, and in the other case
$Q\leq\delta^{-c_*}$ and $L\leq\delta^{-c_*/6}$.
If \eqref{gi-shell-bound} fails with a sufficiently large implied
constant, Proposition~\ref{gi-extraction} applies with
$\epsilon$ a fixed large multiple of $QL^6\delta^{c_*}$; if this
exceeds the permitted range, the elementary bound already proves the
claim.  The extracted density and normalized amplitude are at least
a fixed positive constant times $\delta^{4c_*}$, hence at least
$\delta^{5c_*}$ after decreasing $\delta_0$.

Set $\eta=\delta^{a_*}$.  If an extracted orbit were totally
$\eta$-equidistributed, Lemma~\ref{lem:gn-norm-transfer}, with the
progression analysis above, would bound the extracted sum by
\[
 O(M\delta^{a_*\kappa}L^e),\qquad e=0\text{ or }1.
\]
Here $M=\lceil N/r\rceil\geq N^{1/3}/2$ in Type I and
$M=B\geq N^{1/3}/4$ in Type II.  The rationality and modulus hypotheses
hold because $c_*<a_*\kappa$.  The lower-parameter hypothesis holds
because $a_*s_*\leq\sigma/12$; increasing the fixed threshold for $N$
absorbs the constants two and four.  Since $a_*\kappa=32c_*$, the
displayed bound is smaller than $M\delta^{5c_*}$, even for $e=1$.
This is a contradiction.  Thus every extracted orbit fails total
$\eta$-equidistribution.  Its density is at least $\eta$, since
$a_*>5c_*$.  Also $N\geq\eta^{-K_{\rm rec}}$ by
$a_*s_*K_{\rm rec}\leq1/4$, and $Q\leq\eta^{-1}$.
Lemma~\ref{gi-recombination} now contradicts total
$\delta$-equidistribution, because
$\eta^{C_{\rm rec}}=\delta^{1/4}>\delta$.  This proves
\eqref{gi-shell-bound}.  Fixed values $\delta\geq\delta_0$ are covered
by the elementary bound with a larger constant.

We spell out the passage through all scales.  Use exact real endpoints
$M_j=2^{-j}N$ and retain shells with $M_j\geq2\sqrt\delta N$.
The discarded norms are at most $2\sqrt\delta N$ and their total
absolute weight is $O(\sqrt\delta NL)$ by Lemma~\ref{gi-moments}.
There is no rounding of ideal-norm cutoffs.  On $[\lfloor M_j\rfloor]$
the orbit is totally $2\sqrt\delta$-equidistributed: every progression
of the requisite length has at least $\delta N$ terms for sufficiently
large $N$, and the original discrepancy bound is stronger.
For $\delta\geq N^{-s_*/4}$ the lower-parameter condition for the shell
bound holds on every retained scale, since
$M_j\geq N^{1-s_*/8}$ up to fixed constants.  The progression is merely
intersected with the shell; its common difference remains at most $Q$.
Summing the shell estimates and using $\sum_jM_j\leq2N$ gives
\begin{equation}\label{gi-full-intermediate}
                  O(QNL^7\delta^{c_*/2}).
\end{equation}
For $1/N<\delta<N^{-s_*/4}$ replace $\delta$ by
$\delta^{s_*/8}\geq N^{-s_*/8}$.  Total equidistribution at the
original parameter implies it at this larger one.  The exponent in
\eqref{gi-full-intermediate} becomes $c_*s_*/16$.
For $\delta\leq1/N$, singleton progressions imply
$|F(g(n)\Gamma)|\leq\delta$ at each integer $n$, so the elementary
$O(NL)$ weight bound suffices.  All remaining bounded $N$ are absorbed
by the fixed constant.  Thus one may take
\begin{equation}\label{gi-final-exponent}
                         c_{\rm G}=c_*s_*/16>0.
\end{equation}
Its dependence is only on $m,d$ and the corresponding norm-transfer
exponents; $H_0$ enters constants and thresholds, not this exponent.
\end{proof}

\begin{corollary}[Ordinary prime term]\label{gi-ordinary-primes}
Under the nilmanifold hypotheses of Proposition~\ref{gi-minor-arc},
there is $c_{\mathbb Q}>0$, depending only on $m,d$, such that
\begin{equation}\label{gi-ordinary-bound}
 \left|\sum_{n\leq N}\Lambda'(n)1_P(n)F(g(n)\Gamma)\right|
 \ll_{m,d} QN\log(2N)^7\delta^{c_{\mathbb Q}}\|F\|_{\mathrm{Lip}},
\end{equation}
where $\Lambda'(p)=\log p$ on rational primes and is zero otherwise.
The same bound holds for ordinary von Mangoldt $\Lambda$.
\end{corollary}
\begin{proof}
For positive integers the convolution identities
$\mu*1=1_{n=1}$, $1*\Lambda=\log$, and $\mu*\log=\Lambda$ give
\eqref{gi-convolution-identity} with exactly the displayed norm cutoffs
replaced by integer cutoffs.  Its coefficients are bounded by
$\sum_{c\mid d}\Lambda(c)=\log d$.  Integer multiplicity is one,
so both moment bounds in Lemma~\ref{gi-moments} remain valid.  The
two Cauchy--Schwarz calculations \eqref{gi-first-cauchy} and
\eqref{gi-second-cauchy} apply to the resulting finite integer sums
and give Proposition~\ref{gi-extraction} with the same conservative
parameter $c_0\epsilon^4/L^{24}$.  Lemma~\ref{gi-recombination} concerns
only integer dilation parameters and is unchanged.

The unweighted replacement for norm transfer follows directly from
total equidistribution.  If an interval progression $J\subseteq[X]$
has $|J|\geq\zeta X$, then
\[
 \left|\sum_{n\in J}F(g(n)\Gamma)\right|
       \leq |J|\zeta\|F\|_{\mathrm{Lip}}
       \leq X\zeta\|F\|_{\mathrm{Lip}}.
\]
If $|J|<\zeta X$, its cardinality and the supremum norm give the same
bound.  This treats every modulus and every truncated interval.
Partial summation supplies the logarithmic weight; the identical
long-progression/short-progression dichotomy treats the mean-zero
product function $F\otimes\overline F$ on the product nilmanifold.
Thus the exponent choices and all shell steps in the proof of
Proposition~\ref{gi-minor-arc} apply, with any sufficiently small
fixed positive transfer exponents, to give \eqref{gi-ordinary-bound}
for $\Lambda$.  No extension of a Mobius theorem to $\Lambda$ is assumed.

Finally proper prime powers contribute at most
$O(\sqrt N\log(2N)^2\|F\|_{\mathrm{Lip}})$, by counting at most
$N^{1/k}$ possible bases for each $2\leq k\leq\log_2N$.
Shrink $c_{\mathbb Q}$ below $1/4$.  If $\delta>1/N$, this error is
absorbed because $N\delta^{c_{\mathbb Q}}\geq N^{1-c_{\mathbb Q}}$.
If $\delta\leq1/N$, singleton equidistribution and
$\sum_{n\leq N}\Lambda'(n)\leq N\log(2N)$ prove the bound directly.
This gives the assertion for $\Lambda'$ as well.
\end{proof}

% Scope: these are mean-zero minor-arc estimates.  Progression main terms,
% fixed-field Hecke Siegel--Walfisz, factorization and the final cube/counting
% argument remain separately proved inputs; no integration is asserted here.
